\section{Results}
\label{sec:results}

We organize the results around three questions. First, does specialist merging
improve reasoning on clean and individually obfuscated code? Second, do the
adaptations learned from individual obfuscations still work when those
obfuscations are stacked? Third, how do the methods behave outside the training
setting, on an unseen obfuscation family and on reverse reasoning?


% ============================================================
% RQ1
% ============================================================

\subsection{RQ1: Does Specialist Merging Improve Program Reasoning Across
Clean and Obfuscated Code?}
\label{sec:rq1}

Table~\ref{tab:rq1_performance} compares performance on clean code and on the
individual obfuscations used to train the specialists.

\begin{table*}[!t]
\centering
\small
\setlength{\tabcolsep}{4.5pt}
\renewcommand{\arraystretch}{1.02}

\caption{\textbf{Performance gains across clean and individually obfuscated
code.}
Forward exact-match accuracy on clean code (\texttt{L0}) and mean accuracy
over the five individual seen obfuscations
(\texttt{L1b}, \texttt{L1r}, \texttt{L2}, \texttt{S1}, and \texttt{S2}).
Green and red denote performance above and below a readable untuned baseline,
respectively. Bold denotes the strongest adapted method in each row.
``--'' indicates that no readable untuned baseline is available.}
\label{tab:rq1_performance}

\begin{tabular}{@{}llrrrrrr@{}}
\toprule
\textbf{Model} &
\textbf{Condition} &
\textbf{Base} &
\textbf{Clean} &
\textbf{Breadth} &
\textbf{KL} &
\textbf{Router} &
\textbf{Merge} \\
\midrule

CodeLlama-7B
& Clean
& .257
& \textcolor{green!55!black}{.427}
& \textcolor{green!55!black}{.413}
& \textcolor{green!55!black}{.422}
& \textbf{\textcolor{green!55!black}{.429}}
& \textcolor{green!55!black}{.426} \\
& Singles
& .194
& \textcolor{green!55!black}{.377}
& \textcolor{green!55!black}{.389}
& \textbf{\textcolor{green!55!black}{.402}}
& \textcolor{green!55!black}{.393}
& \textcolor{green!55!black}{.380} \\

\addlinespace[2pt]

CodeLlama-13B
& Clean
& .252
& \textcolor{green!55!black}{.469}
& \textcolor{green!55!black}{.446}
& \textcolor{green!55!black}{.462}
& \textbf{\textcolor{green!55!black}{.477}}
& \textbf{\textcolor{green!55!black}{.477}} \\
& Singles
& .215
& \textcolor{green!55!black}{.410}
& \textcolor{green!55!black}{.423}
& \textbf{\textcolor{green!55!black}{.442}}
& \textcolor{green!55!black}{.433}
& \textcolor{green!55!black}{.424} \\

\addlinespace[2pt]

CodeLlama-34B
& Clean
& --
& .520
& .494
& .525
& .526
& \textbf{.527} \\
& Singles
& .214
& \textcolor{green!55!black}{.462}
& \textcolor{green!55!black}{.469}
& \textbf{\textcolor{green!55!black}{.496}}
& \textcolor{green!55!black}{.484}
& \textcolor{green!55!black}{.460} \\

\addlinespace[2pt]

Llama-3.1-8B
& Clean
& .258
& \textcolor{green!55!black}{.448}
& \textcolor{green!55!black}{.426}
& \textcolor{green!55!black}{.445}
& \textbf{\textcolor{green!55!black}{.475}}
& \textcolor{green!55!black}{.462} \\
& Singles
& .219
& \textcolor{green!55!black}{.390}
& \textcolor{green!55!black}{.393}
& \textcolor{green!55!black}{.416}
& \textbf{\textcolor{green!55!black}{.427}}
& \textcolor{green!55!black}{.406} \\

\addlinespace[2pt]

StarCoder2-15B
& Clean
& --
& .542
& .529
& .532
& \textbf{.546}
& .539 \\
& Singles
& --
& .482
& .492
& \textbf{.512}
& .503
& .475 \\

\addlinespace[2pt]

Gemma-3-12B
& Clean
& --
& .520
& .507
& .502
& \textbf{.530}
& .523 \\
& Singles
& --
& .462
& .474
& .485
& \textbf{.492}
& .479 \\

\addlinespace[2pt]

CodeGemma-7B
& Clean
& .237
& \textcolor{green!55!black}{.465}
& \textcolor{green!55!black}{.439}
& \textcolor{green!55!black}{.445}
& \textbf{\textcolor{green!55!black}{.474}}
& \textcolor{green!55!black}{.465} \\
& Singles
& .200
& \textcolor{green!55!black}{.410}
& \textcolor{green!55!black}{.422}
& \textcolor{green!55!black}{.429}
& \textbf{\textcolor{green!55!black}{.449}}
& \textcolor{green!55!black}{.414} \\

\addlinespace[2pt]

Granite-3.1-8B
& Clean
& .285
& \textcolor{green!55!black}{.430}
& \textcolor{green!55!black}{.416}
& \textcolor{green!55!black}{.396}
& \textbf{\textcolor{green!55!black}{.438}}
& \textcolor{green!55!black}{.433} \\
& Singles
& .262
& \textcolor{green!55!black}{.363}
& \textcolor{green!55!black}{.391}
& \textcolor{green!55!black}{.378}
& \textbf{\textcolor{green!55!black}{.402}}
& \textcolor{green!55!black}{.377} \\

\bottomrule
\end{tabular}
\end{table*}

\paragraph{Merging improves clean-code performance.}
On \texttt{L0}, the merged adapter substantially exceeds the untuned model
wherever a readable baseline is available. It also matches or exceeds direct
clean-code tuning on six of eight models and trails by only 0.1 and 0.3
percentage points on the remaining two. Combining several specialists
therefore does not require giving up the gain obtained from clean-code
training.

Because the default merge contains the \texttt{L0} specialist, this does not
show that the obfuscation specialists alone cause the clean-code improvement.
Rather, the improvement learned by the clean specialist remains intact after
the other specialists are incorporated.

\paragraph{The improvement extends to individual obfuscations.}
The same merged adapter is also substantially stronger than the untuned model
on the individual identifier and structural obfuscations wherever a readable
baseline is available. KL-consistency and routing often achieve higher peak
accuracy on these individual conditions, but the merge combines improvements
across clean and obfuscated code within a single adapter.

\begin{takeaway}
\textbf{Takeaway.}
Specialist merging produces large gains over the untuned model on both clean
and individually obfuscated code. On clean code, it also matches or exceeds
clean-only tuning on six of eight models, showing that multiple specialized
adaptations can be combined without giving up the gains of direct clean-code
training.
\end{takeaway}

\FloatBarrier


% ============================================================
% RQ2
% ============================================================

\subsection{RQ2: Do Separately Learned Adaptations Work When Obfuscations
Are Stacked?}
\label{sec:rq2}

Table~\ref{tab:rq2_depth} evaluates the main compositional setting of the
paper. Each system is trained only on individual transformations, while
evaluation applies two, three, or four obfuscations to the same program.
No system sees these exact stacks during training. The table also includes
each untuned model's performance on clean \texttt{L0} code as a reference.
\begin{table*}[!t]
\centering
\scriptsize
\setlength{\tabcolsep}{2.7pt}
\renewcommand{\arraystretch}{0.95}

\caption{\textbf{Stacked-obfuscation performance relative to clean code.}
Forward exact-match accuracy by stack depth. \textbf{Base \texttt{L0}} is the
untuned model on clean code; $\Delta$ vs.\ \texttt{L0} is
\textbf{Merge $-$ Base \texttt{L0}} in percentage points. Positive values mean
that the merged model on obfuscated code outperforms the untuned model on clean
code. }
\label{tab:rq2_depth}

\begin{tabular}{@{}lrrrrrrrrr@{}}
\toprule
\textbf{Model} &
\textbf{Base \texttt{L0}} &
\textbf{Depth} &
\textbf{Stack base} &
\textbf{Clean} &
\textbf{Breadth} &
\textbf{KL} &
\textbf{Router} &
\textbf{Merge} &
\textbf{$\Delta$ vs.\ \texttt{L0} (pp)} \\
\midrule

CodeLlama-7B
& .257
& 2
& .157
& \textcolor{green!55!black}{.327}
& \textcolor{green!55!black}{.365}
& \textbf{\textcolor{green!55!black}{.378}}
& \textcolor{green!55!black}{.363}
& \textcolor{green!55!black}{.339}
& \textcolor{green!55!black}{+8.2} \\
&
& 3
& .135
& \textcolor{green!55!black}{.328}
& \textcolor{green!55!black}{.366}
& \textbf{\textcolor{green!55!black}{.373}}
& \textcolor{green!55!black}{.359}
& \textcolor{green!55!black}{.328}
& \textcolor{green!55!black}{+7.1} \\
&
& 4
& .089
& \textcolor{green!55!black}{.278}
& \textcolor{green!55!black}{.336}
& \textbf{\textcolor{green!55!black}{.356}}
& \textcolor{green!55!black}{.334}
& \textcolor{green!55!black}{.290}
& \textcolor{green!55!black}{+3.3} \\

\addlinespace[2pt]

CodeLlama-13B
& .252
& 2
& .185
& \textcolor{green!55!black}{.362}
& \textcolor{green!55!black}{.394}
& \textbf{\textcolor{green!55!black}{.407}}
& \textcolor{green!55!black}{.404}
& \textcolor{green!55!black}{.381}
& \textcolor{green!55!black}{+12.9} \\
&
& 3
& .159
& \textcolor{green!55!black}{.354}
& \textcolor{green!55!black}{.395}
& \textbf{\textcolor{green!55!black}{.406}}
& \textcolor{green!55!black}{.399}
& \textcolor{green!55!black}{.375}
& \textcolor{green!55!black}{+12.3} \\
&
& 4
& .117
& \textcolor{green!55!black}{.300}
& \textcolor{green!55!black}{.353}
& \textcolor{green!55!black}{.360}
& \textbf{\textcolor{green!55!black}{.363}}
& \textcolor{green!55!black}{.320}
& \textcolor{green!55!black}{+6.8} \\

\addlinespace[2pt]

CodeLlama-34B
& --
& 2
& .180
& \textcolor{green!55!black}{.401}
& \textcolor{green!55!black}{.432}
& \textbf{\textcolor{green!55!black}{.459}}
& \textcolor{green!55!black}{.437}
& \textcolor{green!55!black}{.408}
& -- \\
&
& 3
& .141
& \textcolor{green!55!black}{.394}
& \textcolor{green!55!black}{.432}
& \textbf{\textcolor{green!55!black}{.452}}
& \textcolor{green!55!black}{.436}
& \textcolor{green!55!black}{.407}
& -- \\
&
& 4
& .105
& \textcolor{green!55!black}{.316}
& \textcolor{green!55!black}{.372}
& \textbf{\textcolor{green!55!black}{.413}}
& \textcolor{green!55!black}{.380}
& \textcolor{green!55!black}{.351}
& -- \\

\addlinespace[2pt]

Llama-3.1-8B
& .258
& 2
& .182
& \textcolor{green!55!black}{.332}
& \textcolor{green!55!black}{.368}
& \textcolor{green!55!black}{.388}
& \textbf{\textcolor{green!55!black}{.394}}
& \textcolor{green!55!black}{.358}
& \textcolor{green!55!black}{+10.0} \\
&
& 3
& .168
& \textcolor{green!55!black}{.322}
& \textcolor{green!55!black}{.364}
& \textcolor{green!55!black}{.380}
& \textbf{\textcolor{green!55!black}{.389}}
& \textcolor{green!55!black}{.355}
& \textcolor{green!55!black}{+9.7} \\
&
& 4
& .140
& \textcolor{green!55!black}{.263}
& \textcolor{green!55!black}{.335}
& \textbf{\textcolor{green!55!black}{.358}}
& \textcolor{green!55!black}{.342}
& \textcolor{green!55!black}{.310}
& \textcolor{green!55!black}{+5.2} \\

\addlinespace[2pt]

StarCoder2-15B
& --
& 2
& --
& .428
& .452
& \textbf{.480}
& .459
& .422
& -- \\
&
& 3
& --
& .431
& .455
& \textbf{.478}
& .456
& .425
& -- \\
&
& 4
& --
& .357
& .412
& \textbf{.437}
& .401
& .358
& -- \\

\addlinespace[2pt]

Gemma-3-12B
& --
& 2
& --
& .395
& .436
& \textbf{.448}
& .440
& .409
& -- \\
&
& 3
& --
& .388
& .432
& \textbf{.449}
& .444
& .410
& -- \\
&
& 4
& --
& .302
& .394
& \textbf{.413}
& .383
& .332
& -- \\

\addlinespace[2pt]

CodeGemma-7B
& .237
& 2
& .164
& \textcolor{green!55!black}{.354}
& \textcolor{green!55!black}{.396}
& \textcolor{green!55!black}{.395}
& \textbf{\textcolor{green!55!black}{.403}}
& \textcolor{green!55!black}{.358}
& \textcolor{green!55!black}{+12.1} \\
&
& 3
& .147
& \textcolor{green!55!black}{.344}
& \textcolor{green!55!black}{.390}
& \textcolor{green!55!black}{.394}
& \textbf{\textcolor{green!55!black}{.408}}
& \textcolor{green!55!black}{.351}
& \textcolor{green!55!black}{+11.4} \\
&
& 4
& .124
& \textcolor{green!55!black}{.271}
& \textbf{\textcolor{green!55!black}{.361}}
& \textcolor{green!55!black}{.359}
& \textcolor{green!55!black}{.360}
& \textcolor{green!55!black}{.303}
& \textcolor{green!55!black}{+6.6} \\

\addlinespace[2pt]

Granite-3.1-8B
& .285
& 2
& .236
& \textcolor{green!55!black}{.314}
& \textbf{\textcolor{green!55!black}{.364}}
& \textcolor{green!55!black}{.356}
& \textcolor{green!55!black}{.363}
& \textcolor{green!55!black}{.330}
& \textcolor{green!55!black}{+4.5} \\
&
& 3
& .240
& \textcolor{green!55!black}{.309}
& \textbf{\textcolor{green!55!black}{.362}}
& \textcolor{green!55!black}{.354}
& \textcolor{green!55!black}{.361}
& \textcolor{green!55!black}{.326}
& \textcolor{green!55!black}{+4.1} \\
&
& 4
& .203
& \textcolor{green!55!black}{.258}
& \textcolor{green!55!black}{.325}
& \textbf{\textcolor{green!55!black}{.334}}
& \textcolor{green!55!black}{.331}
& \textcolor{green!55!black}{.281}
& \textcolor{red!70!black}{-0.4} \\

\bottomrule
\end{tabular}
\end{table*}

\paragraph{Training on individual obfuscations transfers to unseen stacks.}
Across all eight models, adaptation remains effective when transformations are
combined. Wherever a readable untuned baseline is available, every adapted
system remains substantially stronger than the base model even at depth 4.
The gain therefore does not depend on seeing the exact stack during training.

\paragraph{Adapted models can outperform the original model even on clean code.}
The clean-code reference in Table~\ref{tab:rq2_depth} provides a stronger
comparison than the stack baseline alone. Although additional obfuscations make
the task progressively harder for the untuned model, adaptation is strong
enough to more than recover that loss.

At depth 4, the strongest adapted method exceeds the original model's
\texttt{L0} accuracy on every model for which that clean baseline is readable,
by 4.9--12.4 percentage points. The merge itself exceeds the clean-code
baseline on four of those five models: by 3.3 points on CodeLlama-7B, 6.8 on
CodeLlama-13B, 5.2 on Llama-3.1-8B, and 6.6 on CodeGemma-7B. On
Granite-3.1-8B, it is only 0.4 points below the clean baseline.

For example, CodeLlama-7B falls from .257 on clean code to .089 when four
obfuscations are stacked. After adaptation, however, all methods score between
.278 and .356 on that same depth-4 condition. In other words, an adapted model
reasoning over heavily obfuscated code can outperform the original model
reasoning over unobfuscated code.

\paragraph{Fine-tuning and routing achieve the highest peak accuracy.}
KL-consistency is strongest in most model--depth comparisons, while routing
and breadth lead the remaining cases. Merging does not maximize forward
accuracy in this setting, but it remains far above the untuned stack baseline
and usually above the untuned clean-code baseline as well.

\begin{takeaway}
\textbf{Takeaway.}
Training on individual obfuscations transfers strongly to unseen stacks.
Even at depth 4, the strongest adapted system outperforms the untuned model
on clean code for every model with a readable \texttt{L0} baseline. The
merged adapter itself does so on four of five such models and is within
0.4 points on the fifth, despite never training directly on stacked examples.
\end{takeaway}

\FloatBarrier


% ============================================================
% RQ3
% ============================================================

\subsection{RQ3: How Well Do the Methods Hold Up Outside the Training Setting?}
\label{sec:rq3}

Table~\ref{tab:rq3_stress} considers two harder cases: an obfuscation family
that is never used for training and a reverse-reasoning task that none of the
methods are fine-tuned to perform.

\begin{table*}[!t]
\centering
\small
\setlength{\tabcolsep}{4pt}
\renewcommand{\arraystretch}{0.98}

\caption{\textbf{Performance beyond the training setting.}
Forward exact-match accuracy over the held-out encoding-family block and
execution-graded accuracy on the reverse-reasoning task.
Green and red denote performance above and below a readable untuned baseline,
respectively. Bold denotes the strongest adapted method in each row.}
\label{tab:rq3_stress}

\begin{tabular}{@{}llrrrrrr@{}}
\toprule
\textbf{Model} &
\textbf{Test} &
\textbf{Base} &
\textbf{Clean} &
\textbf{Breadth} &
\textbf{KL} &
\textbf{Router} &
\textbf{Merge} \\
\midrule

CodeLlama-7B
& Held-out
& .120
& \textcolor{green!55!black}{.259}
& \textcolor{green!55!black}{.243}
& \textbf{\textcolor{green!55!black}{.280}}
& \textcolor{green!55!black}{.272}
& \textcolor{green!55!black}{.266} \\
& Reverse
& .282
& \textcolor{red!70!black}{.266}
& \textcolor{red!70!black}{.249}
& \textcolor{green!55!black}{.287}
& \textcolor{green!55!black}{.287}
& \textbf{\textcolor{green!55!black}{.311}} \\

\addlinespace[2pt]

CodeLlama-13B
& Held-out
& .145
& \textcolor{green!55!black}{.281}
& \textcolor{green!55!black}{.271}
& \textcolor{green!55!black}{.297}
& \textbf{\textcolor{green!55!black}{.305}}
& \textcolor{green!55!black}{.295} \\
& Reverse
& .294
& \textcolor{green!55!black}{.312}
& --
& \textcolor{red!70!black}{.223}
& .294
& \textbf{\textcolor{green!55!black}{.337}} \\

\addlinespace[2pt]

CodeLlama-34B
& Held-out
& .139
& \textcolor{green!55!black}{.307}
& \textcolor{green!55!black}{.302}
& \textbf{\textcolor{green!55!black}{.337}}
& \textcolor{green!55!black}{.330}
& \textcolor{green!55!black}{.314} \\
& Reverse
& .319
& \textcolor{red!70!black}{.295}
& --
& \textcolor{red!70!black}{.260}
& --
& \textbf{\textcolor{green!55!black}{.381}} \\

\addlinespace[2pt]

Llama-3.1-8B
& Held-out
& .135
& \textcolor{green!55!black}{.243}
& \textcolor{green!55!black}{.233}
& \textcolor{green!55!black}{.260}
& \textbf{\textcolor{green!55!black}{.279}}
& \textcolor{green!55!black}{.263} \\
& Reverse
& .322
& \textcolor{red!70!black}{.282}
& \textcolor{red!70!black}{.264}
& \textcolor{red!70!black}{.273}
& \textcolor{red!70!black}{.296}
& \textbf{\textcolor{green!55!black}{.361}} \\

\addlinespace[2pt]

StarCoder2-15B
& Held-out
& --
& .339
& .332
& .349
& \textbf{.358}
& .344 \\
& Reverse
& .361
& \textcolor{red!70!black}{.335}
& \textcolor{red!70!black}{.295}
& \textcolor{red!70!black}{.336}
& \textcolor{red!70!black}{.350}
& \textbf{\textcolor{green!55!black}{.376}} \\

\addlinespace[2pt]

Gemma-3-12B
& Held-out
& .170
& \textcolor{green!55!black}{.290}
& \textcolor{green!55!black}{.276}
& \textcolor{green!55!black}{.296}
& \textbf{\textcolor{green!55!black}{.309}}
& \textcolor{green!55!black}{.295} \\
& Reverse
& .372
& \textcolor{red!70!black}{.345}
& \textcolor{green!55!black}{.385}
& \textcolor{red!70!black}{.355}
& \textbf{\textcolor{green!55!black}{.406}}
& \textcolor{green!55!black}{.403} \\

\addlinespace[2pt]

CodeGemma-7B
& Held-out
& .149
& \textcolor{green!55!black}{.266}
& \textcolor{green!55!black}{.271}
& \textcolor{green!55!black}{.262}
& \textbf{\textcolor{green!55!black}{.294}}
& \textcolor{green!55!black}{.267} \\
& Reverse
& .218
& \textbf{\textcolor{green!55!black}{.303}}
& \textcolor{green!55!black}{.284}
& --
& \textcolor{green!55!black}{.284}
& \textcolor{green!55!black}{.287} \\

\addlinespace[2pt]

Granite-3.1-8B
& Held-out
& .168
& \textcolor{green!55!black}{.218}
& \textcolor{green!55!black}{.214}
& \textcolor{green!55!black}{.204}
& \textbf{\textcolor{green!55!black}{.256}}
& \textcolor{green!55!black}{.234} \\
& Reverse
& .285
& \textcolor{red!70!black}{.202}
& \textcolor{red!70!black}{.217}
& \textcolor{red!70!black}{.150}
& \textcolor{red!70!black}{.247}
& \textbf{\textcolor{red!70!black}{.270}} \\

\bottomrule
\end{tabular}
\end{table*}

\paragraph{Routing performs best on the unseen obfuscation family.}
Routing achieves the highest forward accuracy on six models, while
KL-consistency is strongest on the remaining two. The router is also above
the merge on all eight models in this block. Merging therefore does not
provide the highest forward accuracy when the transformation itself is new.

Breadth and merge provide a more direct comparison because they use the same
six training conditions. Across the full held-out block, merging exceeds
breadth on seven of eight models. The router, by contrast, keeps a larger bank
of specialists and chooses among them at inference time, whereas the merge
combines its specialists beforehand and deploys only one adapter.

\paragraph{Standard fine-tuning degrades reverse reasoning.}
The reverse task asks the model to find an input that produces a given output.
None of the adaptation methods are trained on this task.

Here, the pattern is almost the opposite of the forward results. Clean-code,
breadth, and KL fine-tuning frequently reduce reverse accuracy below the
untuned model. Merging instead improves over the untuned model on seven of
eight models by 3.0--6.4 percentage points and is the strongest adapted method
on six of eight models.

These results suggest that the SFT-based methods often learn the trained
input--output mapping too narrowly. Their strong forward performance does not
consistently transfer when the same program must be reasoned about in the
opposite direction. In inspection of the reverse-task outputs, the fine-tuned
models frequently return the program's \emph{forward output} even though the
prompt asks for an input that would produce a specified value. This behavior
is consistent with task-specific memorization of the forward prediction
pattern rather than learning a more general relation between program inputs
and outputs.

Merging behaves differently. The specialists are optimized independently and
combined only after training, rather than being jointly optimized toward the
same forward objective. The merged model gives up some peak forward accuracy,
but retains substantially more ability to answer a different question about
the same program.

Reverse accuracy is graded by execution because more than one input may
produce the requested output. Any generated input that returns the correct
value is counted as correct. Approximately 15\% of items have a trivial target
value that a well-formed guess can recover, so
Appendix~\ref{app:absmaster} also reports exact reference-input recovery.

\begin{takeaway}
\textbf{Takeaway.}
Fine-tuning and routing provide the strongest peak forward performance, but
that specialization does not reliably transfer beyond the trained task.
The SFT baselines often lose reverse-reasoning performance and frequently
fall back to predicting the forward output when asked to reason in reverse.
Merging sacrifices some peak forward accuracy but retains substantially more
reverse-reasoning ability.
\end{takeaway}

\FloatBarrier